\documentclass[11pt,a4paper]{article}
\usepackage{fontspec,amsmath,amssymb,amsthm,microtype}
\setmainfont{lmroman10-regular.otf}[BoldFont=lmroman10-bold.otf,ItalicFont=lmroman10-italic.otf,BoldItalicFont=lmroman10-bolditalic.otf]
\usepackage{xcolor}
\usepackage[most]{tcolorbox}
\usepackage{mathrsfs}
\providecommand{\lt}{<}
\providecommand{\gt}{>}
\setmonofont{DejaVuSansMono.ttf}[BoldFont=DejaVuSansMono-Bold.ttf,ItalicFont=DejaVuSansMono-Oblique.ttf,BoldItalicFont=DejaVuSansMono-BoldOblique.ttf,Scale=0.86]
\newfontfamily\leanfallback{STIXTwoMath-Regular.otf}[Scale=0.86]
\definecolor{leanrule}{gray}{0.62}
\definecolor{leanback}{gray}{0.962}
\definecolor{leantitle}{gray}{0.22}
\definecolor{leanmuted}{gray}{0.38}
\newtcolorbox{leanbox}[1]{enhanced,breakable,colback=leanback,colframe=leanrule,
  boxrule=0.4pt,arc=0.6mm,left=2.4mm,right=2.4mm,top=1.8mm,bottom=1.8mm,
  before skip=0.8em,after skip=1em,fonttitle=\small\bfseries,coltitle=leantitle,
  colbacktitle=leanback,titlerule=0.3pt,title={#1}}
\newlength{\leancharwidth}
\newenvironment{leancode}{\par\smallskip\begingroup\ttfamily\footnotesize\raggedright
  \settowidth{\leancharwidth}{M}\setlength{\parskip}{0pt}\setlength{\parindent}{0pt}
  \frenchspacing}{\par\endgroup\smallskip}
\newcommand{\leanline}[2]{\hangindent=\dimexpr#1\leancharwidth+4\leancharwidth\relax
  \hangafter=1\noindent\hspace*{#1\leancharwidth}#2\par}
\newcommand{\leansource}[1]{\par{\footnotesize\color{leanmuted}#1}\par}
\newcommand{\leanlabel}[1]{\par\smallskip\noindent{\small\bfseries #1}\par}
\usepackage[margin=25mm]{geometry}
\usepackage[unicode,hidelinks]{hyperref}
\hypersetup{pdfauthor={LangevinLean Project}}
\setlength{\emergencystretch}{3em}
\setlength{\parskip}{0.4em}
\setlength{\parindent}{0pt}
\allowdisplaybreaks[2]
\raggedbottom
\title{\texorpdfstring{Stationary independent inputs for the actual future dynamics}{Stationary independent inputs for the actual future dynamics}}
\author{LangevinLean Project}
\date{}
\begin{document}\maketitle
Draft mathematical exposition; these notes may contain mistakes.


At each deterministic time, the actual state has the Gibbs law and is measurable from the past. Its pair with the whole independent future Brownian path therefore has a product law, and this pair law is independent of the cut time.


\section{\texorpdfstring{The state and future form an actual product input}{The state and future form an actual product input}}


\textbf{Claim.} At \(a\ge0\), the state \(Y_a\) is past-measurable, is independent of \(W_a\), and

\[\operatorname{Law}(Y_a,W_a)=\pi\otimes\operatorname{Law}(W_a).\]

\textbf{Proof.} Apply natural-past adaptation to a finite horizon ending at \(a\), then identify the endpoint with the global solution. Independence of the whole future from the past implies independence from this measurable state. The pair-law factorization follows from independence and the already proved Gibbs stationarity.

\begin{leanbox}{Lean: \texttt{actualState\_\allowbreak{}adapted\_\allowbreak{}pastAt}, \texttt{actualState\_\allowbreak{}future\_\allowbreak{}independent}, \texttt{actualState\_\allowbreak{}future\_\allowbreak{}joint\_\allowbreak{}law}}
\leansource{\texttt{actualState\_\allowbreak{}adapted\_\allowbreak{}pastAt} \quad WassersteinMixing/\allowbreak{}ActualMarkovInput.lean, lines 19–29}
\begin{leancode}
\leanline{0}{theorem\ actualState\_adapted\_pastAt\ (a\ :\ ℝ≥0)\ :}
\leanline{4}{AEStronglyMeasurable[ULA.lpLangevinPast\ (actualProbability\ P\ V)}
\leanline{6}{(CanonicalSemigroup.initial\ P\ V\ hV)\ (actualBrownianNoise\ P\ V\ hV)\ (a\ :\ ℝ)]}
\leanline{6}{(CanonicalSemigroup.state\ P\ V\ hV\ a)\ (actualProbability\ P\ V)}
\end{leancode}
\leansource{\texttt{actualState\_\allowbreak{}future\_\allowbreak{}independent} \quad WassersteinMixing/\allowbreak{}ActualMarkovInput.lean, lines 32–41}
\begin{leancode}
\leanline{0}{theorem\ actualState\_future\_independent\ (a\ :\ ℝ≥0)\ :}
\leanline{4}{IndepFun\ (CanonicalSemigroup.state\ P\ V\ hV\ a)\ (actualBrownianFuture\ P\ a)\ (actualProbability\ P\ V)}
\end{leancode}
\leansource{\texttt{actualState\_\allowbreak{}future\_\allowbreak{}joint\_\allowbreak{}law} \quad WassersteinMixing/\allowbreak{}ActualMarkovInput.lean, lines 44–51}
\begin{leancode}
\leanline{0}{theorem\ actualState\_future\_joint\_law\ (a\ :\ ℝ≥0)\ :}
\leanline{4}{(actualProbability\ P\ V).map\ (fun\ ω\ ↦}
\leanline{6}{(CanonicalSemigroup.state\ P\ V\ hV\ a\ ω,\ actualBrownianFuture\ P\ a\ ω))\ =}
\leanline{6}{(ULA.gibbsMeasure\ V).prod\ ((actualProbability\ P\ V).map\ (actualBrownianFuture\ P\ a))}
\end{leancode}
\leanlabel{How the Lean proof follows the argument}
\texttt{actual\allowbreak{}State\_\allowbreak{}adapted\_\allowbreak{}past\allowbreak{}At} applies the actual integral-solution adaptation theorem. \texttt{actual\allowbreak{}State\_\allowbreak{}future\_\allowbreak{}independent} descends future/\allowbreak{}past independence to the state. \texttt{actual\allowbreak{}State\_\allowbreak{}future\_\allowbreak{}joint\_\allowbreak{}law} uses the exact product pushforward identity.
\end{leanbox}


Earlier results used here: \cite[Section~3.25.2]{LangevinLean2026ULASC}; \cite[Section~3.25.3]{LangevinLean2026ULASC}; \cite[Section~3.26.3]{LangevinLean2026ULASC}; \cite[Section~3.54.4]{LangevinLean2026ULASC}; \cite[Section~3.54.8]{LangevinLean2026ULASC}.

\section{\texorpdfstring{Equality of whole future-input laws}{Equality of whole future-input laws}}


\textbf{Claim.} The shifted Brownian path has the same law at any two cut times, so \(\operatorname{Law}(Y_a,W_a)=\operatorname{Law}(Y_b,W_b)\).



\textbf{Proof.} Finite collections of shifted Brownian coordinates have the same centered Gaussian law, with covariance determined solely by their elapsed times. Equality on finite coordinate cylinders gives equality of the full path laws. The actual vector source is a product of independent coordinate sources; equality therefore transfers to the vector path. Substitute this equality into the preceding product-input factorization.

\begin{leanbox}{Lean: \texttt{preBrownian\_\allowbreak{}shift\_\allowbreak{}path\_\allowbreak{}law}, \texttt{actualBrownianFuture\_\allowbreak{}law\_\allowbreak{}eq}, \texttt{actualState\_\allowbreak{}future\_\allowbreak{}joint\_\allowbreak{}law\_\allowbreak{}eq}}
\leansource{\texttt{preBrownian\_\allowbreak{}shift\_\allowbreak{}path\_\allowbreak{}law} \quad WassersteinMixing/\allowbreak{}ActualMarkovInput.lean, lines 53–63}
\begin{leancode}
\leanline{0}{private\ theorem\ preBrownian\_shift\_path\_law\ \{Ω\ :\ Type*\}\ [MeasurableSpace\ Ω]}
\leanline{4}{(Q\ :\ Measure\ Ω)\ [IsProbabilityMeasure\ Q]\ (B\ :\ ℝ≥0\ →\ Ω\ →\ ℝ)}
\leanline{4}{(hB\ :\ IsPreBrownianReal\ B\ Q)\ (hBmeas\ :\ ∀\ t,\ Measurable\ (B\ t))\ (a\ b\ :\ ℝ≥0)\ :}
\leanline{4}{Q.map\ (fun\ ω\ t\ ↦\ B\ (a\ +\ t)\ ω\ -\ B\ a\ ω)\ =\ Q.map\ (fun\ ω\ t\ ↦\ B\ (b\ +\ t)\ ω\ -\ B\ b\ ω)}
\end{leancode}
\leansource{\texttt{actualBrownianFuture\_\allowbreak{}law\_\allowbreak{}eq} \quad WassersteinMixing/\allowbreak{}ActualMarkovInput.lean, lines 66–98}
\begin{leancode}
\leanline{0}{theorem\ actualBrownianFuture\_law\_eq\ (a\ b\ :\ ℝ≥0)\ :}
\leanline{4}{(actualProbability\ P\ V).map\ (actualBrownianFuture\ P\ a)\ =}
\leanline{6}{(actualProbability\ P\ V).map\ (actualBrownianFuture\ P\ b)}
\end{leancode}
\leansource{\texttt{actualState\_\allowbreak{}future\_\allowbreak{}joint\_\allowbreak{}law\_\allowbreak{}eq} \quad WassersteinMixing/\allowbreak{}ActualMarkovInput.lean, lines 101–107}
\begin{leancode}
\leanline{0}{theorem\ actualState\_future\_joint\_law\_eq\ (a\ b\ :\ ℝ≥0)\ :}
\leanline{4}{(actualProbability\ P\ V).map\ (fun\ ω\ ↦}
\leanline{6}{(CanonicalSemigroup.state\ P\ V\ hV\ a\ ω,\ actualBrownianFuture\ P\ a\ ω))\ =}
\leanline{4}{(actualProbability\ P\ V).map\ (fun\ ω\ ↦}
\leanline{6}{(CanonicalSemigroup.state\ P\ V\ hV\ b\ ω,\ actualBrownianFuture\ P\ b\ ω))}
\end{leancode}
\leanlabel{How the Lean proof follows the argument}
\texttt{pre\allowbreak{}Brownian\_\allowbreak{}shift\_\allowbreak{}path\_\allowbreak{}law} proves equality of path measures from finite-dimensional Gaussian data. \texttt{actual\allowbreak{}Brownian\allowbreak{}Future\_\allowbreak{}law\_\allowbreak{}eq} applies it to the actual coordinate product. The final theorem rewrites both exact product input laws.
\end{leanbox}


Earlier results used here: \cite[Section~3.22.1]{LangevinLean2026ULASC}; \cite[Section~3.8.11]{LangevinLean2026ULASC}; \cite[Section~3.8.16]{LangevinLean2026ULASC}; \cite[Section~3.8.17]{LangevinLean2026ULASC}; \cite[Section~3.8.19]{LangevinLean2026ULASC}; \cite[Section~3.8.20]{LangevinLean2026ULASC}; \cite[Section~3.8.21]{LangevinLean2026ULASC}.

\begin{thebibliography}{ULA-LSI}
\bibitem[ULA-SC]{LangevinLean2026ULASC}
LangevinLean Project. \emph{Unadjusted Langevin under Strong Convexity}. LL-01, 2026. Draft report, content revision d34ead06987c. \url{https://odakyildiz.github.io/LangevinLean/docs/pdf/UnadjustedLangevin/FullProof.pdf}.
\end{thebibliography}

\end{document}
